% =========================================================================
% Capítulo 1: Cinemática
% Tema 05: Problemas de mezcla compleja
% Problema 02: Barra en aro deslizante
% =========================================================================
\subsubsection{Barra en aro deslizante}

Una barra delgada de longitud $l$ está articulada a un aro como se muestra en la Figura \ref{fig:barra_deslizante}. El aro se desliza libremente a lo largo de una pista horizontal fija. Si se sabe que $x$ es el desplazamiento horizontal del aro y que $\theta$ describe la orientación de la barra en relación con la dirección vertical, determine las siguientes magnitudes según las ve un observador fijo en la pista (marco de referencia inercial):

\begin{enumerate}
    \item El vector velocidad del extremo libre de la barra (punto $P$).
    \item El vector aceleración del extremo libre de la barra (punto $P$).
\end{enumerate}

\begin{figure}[!ht]
    \centering
    \begin{tikzpicture}[>=Latex, scale=1.5]
        % --- Parámetros ---
        \def\L{2.5}         % Longitud de la barra l
        \def\ThetaAng{45}   % Ángulo theta
        \def\TrackW{5.0}    % Ancho de la pista

        % --- 1. Pista Horizontal (Fija) ---
        % Sombreado de los extremos fijos
        \fill[pattern=north east lines, pattern color=gray!60] (-2.5, -0.3) rectangle (-2.3, 0.3);
        \draw[-,thick] (-2.3,-0.3)--++(0,0.6);
        \fill[pattern=north east lines, pattern color=gray!60] (2.3, -0.3) rectangle (2.5, 0.3);
        \draw[-,thick] (2.3,-0.3)--++(0,0.6);
        
        % Cuerpo de la pista (dos líneas paralelas para dar volumen)
        \draw[thick, fill=gray!10] (-2.3, -0.08) rectangle (2.3, 0.08);

        % --- 2. Aro Deslizante (Punto O) ---
        \coordinate (O) at (0,0);
        \shadedraw[inner color=white, outer color=gray!60, rounded corners=2pt] (-0.2,-0.15) rectangle (0.2,0.15);
        \filldraw (O) circle (1.5pt) node[below left, yshift=-2pt] {$O$};

        % --- 3. Barra delgada (Punto O a P) ---
        \coordinate (P) at ({\L*sin(\ThetaAng)}, {\L*cos(\ThetaAng)});
        
        % Dibujo de la barra con grosor y degradado
        \draw[line width=5pt, line cap=round, black] (O) -- (P);
        \draw[line width=4pt, line cap=round, gray!30] (O) -- (P);
        \node[below right] at ($(O)!0.5!(P)$) {$l$};
        \filldraw (P) circle (0.8pt) node[above right] {$P$};

        % --- 4. Referencias y Ángulo ---
        % Eje vertical de referencia
        \draw[dashed, thick] (O) -- (0, 1.5);
        
        % Arco para el ángulo theta
        \draw[->] (0, 1.0) arc (90:{90-\ThetaAng}:1.0);
        \node at (70:1.2) {$\theta$};

        % --- 5. Indicación de desplazamiento x ---
        \draw[|->, thick] (0, -0.5) -- (1.0, -0.5) node[right] {$x$};

    \end{tikzpicture}
    \caption{Barra de longitud $l$ articulada a un aro que se desplaza horizontalmente una distancia $x$.}
    \label{fig:barra_deslizante}
\end{figure}5

% --- SOLUCIÓN ---------------------------------------------------------------
\begin{solucion}{Barra en aro deslizante (cinemática)}
\textbf{Planteamiento}\\
Hay dos grados de libertad, $x$ y $\theta$. Según la figura la barra está \emph{sobre} el aro y $\theta$ se mide desde la vertical hacia arriba, creciendo en sentido horario: $\theta=0$ es la posición invertida y $\theta=\pi$ la barra colgando.

\patoojo[squareglasses,shovel]{¡Ojo con la figura! La barra apunta hacia \emph{arriba}: $\theta=0$ es la posición invertida (inestable) y $\theta=\pi$ es la barra colgando. Además $\theta$ crece en sentido \emph{horario}, así que $\velang{N}{B}=-\dot\theta\vec{E}_z$.}

\textbf{Marcos de referencia y Sistemas de Coordenadas}
\begin{itemize}
    \item $\mathcal{N}$ La pista (inercial) $\longrightarrow\{\vec{E}_x;\vec{E}_y;\vec{E}_z\}$
    \item $\mathcal{B}$ La barra $\longrightarrow\{\vec{e}_r;\vec{e}_\theta;\vec{E}_z\}$, con $\vec{e}_r$ a lo largo de la barra
\end{itemize}
\begin{equation*}
\begin{split}
    \vec{e}_r&=\sin\theta\,\vec{E}_x+\cos\theta\,\vec{E}_y\\
    \vec{e}_\theta&=\cos\theta\,\vec{E}_x-\sin\theta\,\vec{E}_y
\end{split}
\end{equation*}

\subsubsection*{Cinemática de la partícula}
\textbf{Posición}\\
Tomando como referencia el aro, con $O=(x,0)$ en la pista, el extremo libre $P$ está a una distancia $l$ de $O$ a lo largo de la barra:
$$\vec{r}_P=x\vec{E}_x+l\vec{e}_r=\left(x+l\sin\theta\right)\vec{E}_x+l\cos\theta\,\vec{E}_y$$
Y la velocidad angular de la barra es $\velang{N}{B}=-\dot\theta\vec{E}_z$ (horaria).

\textbf{Velocidad}\\
$$\marc{\vec{v}_P}{N}=\der{\vec{r}_P}{N}=\dot{x}\vec{E}_x+\underbrace{\der{(l\vec{e}_r)}{B}{}}_{\vec{0}}+\velang{N}{B}\times l\vec{e}_r=\dot{x}\vec{E}_x+l\dot\theta\vec{e}_\theta$$
\respuesta{$\displaystyle\marc{\vec{v}_P}{N}=\left(\dot{x}+l\dot\theta\cos\theta\right)\vec{E}_x-l\dot\theta\sin\theta\,\vec{E}_y$}

\textbf{Aceleración}\\
Derivando de nuevo con el teorema del transporte, ahora sobre $\dot{x}\vec{E}_x+l\dot\theta\vec{e}_\theta$:
$$\marc{\vec{a}_P}{N}=\ddot{x}\vec{E}_x+l\ddot\theta\vec{e}_\theta+\velang{N}{B}\times\left(l\dot\theta\vec{e}_\theta\right)=\ddot{x}\vec{E}_x+l\ddot\theta\vec{e}_\theta-l\dot\theta^2\vec{e}_r$$
\respuesta{$\displaystyle\marc{\vec{a}_P}{N}=\left[\ddot{x}+l\left(\ddot\theta\cos\theta-\dot\theta^2\sin\theta\right)\right]\vec{E}_x-l\left(\ddot\theta\sin\theta+\dot\theta^2\cos\theta\right)\vec{E}_y$}

\patotip{Si en cambio piden el centro de masa $G$ de la barra, el cálculo es idéntico reemplazando $l$ por $l/2$ (está a mitad de la barra).}
\end{solucion}
